\documentclass[11pt,a4paper]{article}

\usepackage[T1]{fontenc}
\usepackage[utf8]{inputenc}
\usepackage[brazil]{babel}
\usepackage[a4paper,margin=2.3cm]{geometry}
\usepackage{lmodern}
\usepackage{amsmath}
\usepackage{xcolor}
\usepackage{graphicx}
\usepackage{listings}
\usepackage{enumitem}
\usepackage{booktabs}
\usepackage{longtable}
\usepackage{fancyhdr}
\usepackage[hidelinks]{hyperref}

\definecolor{codebg}{RGB}{245,245,245}
\definecolor{codekw}{RGB}{0,0,160}
\definecolor{codecm}{RGB}{90,120,90}
\definecolor{codestr}{RGB}{160,30,30}
\definecolor{boxbg}{RGB}{235,242,250}

\lstset{
  basicstyle=\ttfamily\footnotesize,
  backgroundcolor=\color{codebg},
  keywordstyle=\color{codekw}\bfseries,
  commentstyle=\color{codecm}\itshape,
  stringstyle=\color{codestr},
  numbers=left, numberstyle=\tiny\color{gray}, numbersep=6pt,
  breaklines=true, breakatwhitespace=false,
  showstringspaces=false, frame=single, rulecolor=\color{gray!40},
  tabsize=2, captionpos=b, columns=fullflexible,
  inputencoding=utf8, extendedchars=true,
  literate={á}{{\'a}}1 {é}{{\'e}}1 {í}{{\'i}}1 {ó}{{\'o}}1 {ú}{{\'u}}1
           {ã}{{\~a}}1 {õ}{{\~o}}1 {â}{{\^a}}1 {ê}{{\^e}}1 {ô}{{\^o}}1
           {ç}{{\c c}}1 {Á}{{\'A}}1 {É}{{\'E}}1 {Í}{{\'I}}1 {Ó}{{\'O}}1
           {Ú}{{\'U}}1 {Ã}{{\~A}}1 {Õ}{{\~O}}1 {Ç}{{\c C}}1
           {à}{{\`a}}1 {—}{{---}}1 {–}{{--}}1 {“}{{``}}1 {”}{{''}}1
           {’}{{'}}1 {€}{{\texteuro}}1 {º}{{\textordmasculine}}1
           {φ}{{$\varphi$}}1 {·}{{$\cdot$}}1
}

\lstdefinestyle{plain}{numbers=none,basicstyle=\ttfamily\small,backgroundcolor=\color{boxbg}}

\pagestyle{fancy}
\fancyhf{}
\lhead{\small Criptoanálise --- AED3}
\rhead{\small Grupo 2}
\cfoot{\thepage}

\title{\textbf{Criptoanálise de um Esquema Híbrido RSA + AES}\\[4pt]
\large Quebra de senha (SHA-256), fatoração RSA via GNFS e recuperação da mensagem}
\author{Grupo 2 --- Gabriel Madureira e Gustavo Brendler}
\date{Algoritmos e Estruturas de Dados III (AED3) --- UFPel\\
Prof. Ulisses Brisolara Correa\\[2pt]\today}

\begin{document}
\maketitle
\thispagestyle{fancy}

\begin{abstract}
\noindent Este relatório documenta, de forma replicável, a criptoanálise completa de um
desafio no tema do \emph{Biuro Szyfrów} (a seção polonesa que quebrou a Enigma). A partir de
quatro artefatos (um \emph{hash} e três arquivos cifrados), recuperamos integralmente a cadeia:
(i)~quebramos o \emph{hash} SHA-256 da senha; (ii)~deciframos a chave pública RSA; (iii)~fatoramos
o módulo RSA de 382~bits com o \emph{General Number Field Sieve} (CADO-NFS); (iv)~reconstruímos a
chave privada e deciframos a chave simétrica; e (v)~deciframos a mensagem final. Registramos
também as abordagens que \textbf{falharam} e sua explicação técnica, conforme exigido.
\end{abstract}

\tableofcontents
\bigskip

\section{Contexto e objetivo}
O enunciado situa-se no tema histórico do \emph{Biuro Szyfrów} e da máquina Enigma. São fornecidos
quatro artefatos:
\begin{itemize}[nosep]
  \item \texttt{key\_for\_rsa\_public.hash} --- \emph{hash} SHA-256 da senha (pista: começa com \texttt{ma});
  \item \texttt{key\_public.en} (157~B) --- chave pública RSA cifrada simetricamente com a senha;
  \item \texttt{key\_for\_message.en} (48~B) --- a chave simétrica da mensagem, cifrada com RSA;
  \item \texttt{message.en} (1482~B) --- a mensagem final, cifrada simetricamente.
\end{itemize}
O objetivo é recuperar a mensagem em claro, documentando todos os procedimentos de modo que o
processo seja reprodutível a partir do zero.

\section{Arquitetura do desafio}
A inspeção dos arquivos (comando \texttt{file} e análise dos cabeçalhos com \texttt{xxd}) revelou a
seguinte cadeia de dependências, típica de um esquema \emph{híbrido} RSA+AES:

\begin{center}
\fbox{\parbox{0.95\linewidth}{\centering
\texttt{senha} $\xrightarrow{\text{SHA-256}}$ \emph{hash} (alvo)\\[2pt]
$\downarrow$ (quebrar o \emph{hash})\\[2pt]
\texttt{senha} $\xrightarrow{\text{AES simétrico}}$ decifra \texttt{key\_public.en} $\Rightarrow$ chave pública RSA $(N,e)$\\[2pt]
$\downarrow$ (fatorar $N$)\\[2pt]
$N=p\cdot q \Rightarrow d=e^{-1}\bmod\varphi(N)$ (chave privada)\\[2pt]
$\downarrow$\\[2pt]
\texttt{key\_for\_message.en} $\xrightarrow{\text{RSA }c^d\bmod N}$ chave simétrica da mensagem\\[2pt]
$\downarrow$\\[2pt]
\texttt{message.en} $\xrightarrow{\text{AES simétrico}}$ \textbf{mensagem em claro}
}}
\end{center}

\noindent Observações de cabeçalho:
\begin{itemize}[nosep]
  \item \texttt{message.en} inicia com \texttt{Salted\_\_} (KDF do \texttt{openssl enc} com \emph{salt});
  \item \texttt{key\_public.en} \emph{não} tem \texttt{Salted\_\_} (cifragem \texttt{-nosalt});
  \item \texttt{key\_for\_message.en} tem exatamente 48~B $=$ 384~bits $\Rightarrow$ ciphertext RSA-384.
\end{itemize}

\section{Etapa 1 --- Quebra do \emph{hash} da senha (SHA-256)}
\label{sec:hash}
O \emph{hash} alvo tem 64 dígitos hexadecimais (256~bits). Como o material da disciplina não ensina
funções de \emph{hash} customizadas, assumimos \textbf{SHA-256 padrão} (e validamos contra outros 6
\emph{hashes} de 256~bits como rede de segurança). A pista é que a senha começa com \texttt{ma}.

\paragraph{Estratégia vencedora.} Após descartar força bruta curta e dicionários simples
(Seção~\ref{sec:falhas}), aplicamos um \textbf{combinador de duas palavras} em C+OpenMP: para cada
palavra iniciada em \texttt{ma} (de um dicionário), concatenamos cada palavra de um dicionário amplo
(com conectores em português opcionais), testando SHA-256 com e sem \texttt{\textbackslash n} final.
A senha foi encontrada:

\begin{center}\fbox{\large\ttfamily senha = marvinminsky}\end{center}

\noindent Trata-se de \textbf{Marvin Minsky}, um dos fundadores da Inteligência Artificial
(\texttt{marvin}+\texttt{minsky}); começa com \texttt{ma}, minúscula, 12 caracteres, \emph{sem}
terminador. Verificação:
\begin{lstlisting}[style=plain,language=bash]
$ echo -n "marvinminsky" | sha256sum
1b78cdb61662342d53da0544a38c87579c4ca0e9b171ad2979f165a6b3e28cce
\end{lstlisting}
que é exatamente o conteúdo de \texttt{key\_for\_rsa\_public.hash}.

\section{Etapa 2 --- Decifração da chave pública RSA}
Com a senha, varremos cifras/modos/KDFs do OpenSSL sobre \texttt{key\_public.en}
(script \texttt{decrypt\_pubkey.py}). O par vencedor foi \textbf{AES-256-CTR, chave derivada por
MD=SHA-256, sem \emph{salt}}:
\begin{lstlisting}[style=plain,language=bash]
$ openssl enc -d -aes-256-ctr -nosalt -md sha256 \
      -pass pass:marvinminsky -in key_public.en
-----BEGIN PUBLIC KEY-----
MEswDQYJKoZIhvcNAQEBBQADOgAwNwIwJb8QgzaQ/v8DNvt/lOiV0PgtAsRh+ZZx
6UVGq9OmaaPvjcJ4p+PmqLHCpYLLVSUbAgMBAAE=
-----END PUBLIC KEY-----
\end{lstlisting}
Extraindo os parâmetros (\texttt{openssl rsa -pubin -text -noout}):
\begin{itemize}[nosep]
  \item Módulo $N$ (382~bits):\\
  \texttt{\footnotesize 0x25bf10833690feff0336fb7f94e895d0f82d02c461f99671e94546abd3a669a3ef8dc278a7e3e6a8b1c2a582cb55251b}
  \item Expoente público $e = 65537$.
\end{itemize}

\section{Etapa 3 --- Fatoração de $N$ (GNFS / CADO-NFS)}
A segurança do RSA repousa na dificuldade de fatorar $N$. Com apenas 382~bits, $N$ é
\emph{fatorável} --- porém \textbf{não} por métodos elementares quando os primos são balanceados
(ver Seção~\ref{sec:falhas}). Empregamos o \emph{General Number Field Sieve} via \textbf{CADO-NFS}:

\begin{lstlisting}[style=plain,language=bash]
$ ./cado-nfs.py 5809694326629984678335471447012067562303627826\
0104987098146141854317060643783262578537565380650848127701261\
61323291 -t 4
\end{lstlisting}

\noindent Resultado (tempo real $\approx$ 2h39; $\approx$ 2h38 de CPU acumulada, 4 \emph{threads}):
\begin{align*}
p &= 1967103988726232522467829353610346625438028501545676588173\\
q &= 2953425116275607543642880749640518195177306268255592720967
\end{align*}
com $p\cdot q = N$ (ambos primos, verificado por Miller-Rabin). O \emph{factordb} não possuía esses
fatores, confirmando que a fatoração foi realizada por nós.

\paragraph{Montagem do \emph{toolchain} sem privilégios de administrador.} O ambiente não possuía
compilador LaTeX, \texttt{pip}, \texttt{cmake}, cabeçalhos do GMP nem ferramentas de fatoração.
Compilamos tudo a partir do código-fonte, em diretório do usuário: \texttt{m4} $\rightarrow$
\texttt{GMP 6.3.0} $\rightarrow$ \texttt{CADO-NFS}, obtendo o \texttt{cmake} via \texttt{pip} (ambiente
virtual) e as dependências Python do CADO (\texttt{flask}, \texttt{requests}) em diretório próprio.

\section{Etapa 4 --- Recuperação da chave simétrica (RSA)}
Com $p,q$ reconstruímos a chave privada: $\varphi(N)=(p-1)(q-1)$ e $d=e^{-1}\bmod\varphi(N)$.
Decifrando $c=$ \texttt{key\_for\_message.en} como inteiro \emph{big-endian}: $m=c^{d}\bmod N$.
O texto recuperado trazia \textbf{padding PKCS\#1 v1.5} (\texttt{00 02 \dots 00 \textbar\ payload}); após
removê-lo, obtivemos a chave simétrica:
\begin{center}\fbox{\large\ttfamily chave da mensagem = shafigoldwas}\end{center}
Referência provável a \textbf{Shafi Goldwasser} (prêmio Turing em criptografia), truncada a 12~bytes
pelo tamanho do bloco RSA-384 --- reforçando o tema ``pioneiros da computação'' também visto na senha.

\section{Etapa 5 --- Decifração da mensagem final}
A mensagem usa \texttt{openssl enc} com \emph{salt}. A combinação vencedora foi
\textbf{AES-256-CFB, MD=SHA-256}:
\begin{lstlisting}[style=plain,language=bash]
$ openssl enc -d -aes-256-cfb -md sha256 \
      -pass pass:shafigoldwas -in message.en
\end{lstlisting}

\noindent Mensagem recuperada (1466~bytes):
\lstinputlisting[style=plain]{../results/MENSAGEM_pkcs1_aes-256-cfb_sha256.txt}

\noindent A mensagem endereça explicitamente o \textbf{Grupo 2} (Gabriel Madureira e Gustavo
Brendler), confirmando o sucesso da criptoanálise.

\section{Abordagens que falharam (e por quê)}
\label{sec:falhas}
Documentamos aqui as principais tentativas malsucedidas e a razão técnica de cada fracasso. (O
registro completo, com mais de 30 variações, está em \texttt{TENTATIVAS.md}.)

\begin{enumerate}[leftmargin=*]
  \item \textbf{Força bruta de senha ``longa'' (15+ caracteres).}
  Guiados por um artigo de segurança indicado em aula (que recomenda \emph{passphrases} de 15+
  caracteres), chegamos a supor uma frase longa. Um espaço de $26^{13}\approx 2.5\times10^{18}$
  (minúsculas, 13 caracteres após \texttt{ma}) a $\sim$27~milhões de \emph{hashes}/s levaria
  $\approx$2900~anos. \emph{Conclusão: inviável} --- e, de fato, a senha era um \textbf{nome de
  12~caracteres}, não uma \emph{passphrase}. Lição: priorizar o \emph{tema} (nomes de pioneiros)
  em vez do tamanho.

  \item \textbf{Ataque ``textbook RSA'' à chave da mensagem (contornar a senha).}
  Tentamos recuperar a chave simétrica diretamente de \texttt{key\_for\_message.en} assumindo
  expoente pequeno ($m=\sqrt[e]{c}$, $e\in\{3,5,7,\dots\}$, \emph{big/little-endian}), sem precisar
  fatorar $N$. Falhou: o expoente real é $e=65537$ e $c\approx 2^{381}$ sofre redução modular por
  $N$, logo \emph{não há raiz inteira exata}. Sem $N$ (portanto sem a senha) não há atalho.

  \item \textbf{Dicionários e \emph{rockyou} + regras (\emph{best64}, \emph{leet}, anos).}
  Testamos $\sim$15 funções de \emph{hash} sobre \emph{rockyou} (filtrado por prefixo \texttt{ma})
  e dicionários PT/EN com \emph{mangling}. Falhou: a senha (\texttt{marvinminsky}) é a
  \emph{concatenação de dois nomes}, não uma palavra isolada nem uma variação trivial --- só um
  \textbf{combinador de 2 palavras} a alcançaria.

  \item \textbf{Fatoração elementar de $N$ (Fermat, Pollard $p\!-\!1$, Brent-$\rho$).}
  Em Python puro, resolvem apenas casos fracos (primos próximos ou com fator pequeno). Para
  $N$ de 382~bits com primos \emph{balanceados} ($\approx$191~bits cada), o custo é proibitivo
  ($\rho$ seria $\sim$$2^{95}$ operações). Por isso foi necessário o \textbf{GNFS (CADO-NFS)}.
\end{enumerate}

\paragraph{Resumo das famílias de ataque testadas.} A Tabela~\ref{tab:resumo} consolida, por
família, as mais de 30 variações registradas em \texttt{TENTATIVAS.md} (mantido como apêndice bruto).
Nenhuma das abordagens de força bruta, dicionário/\emph{rockyou} ou \emph{passphrase} teve êxito;
somente o \textbf{combinador de duas palavras} recuperou a senha.

\begin{table}[h]
\centering\small
\begin{tabular}{p{0.44\linewidth} p{0.28\linewidth} l}
\toprule
\textbf{Família de ataque} & \textbf{Espaço / ferramentas} & \textbf{Resultado}\\
\midrule
Força bruta de prefixo \texttt{ma} ([a-z], [a-z0-9], caixa mista, símbolos; até 8--10 chars) & bilhões de \emph{hashes}; 7 funções de 256 bits (C+OpenMP) & falhou\\
Dicionários $+$ \emph{rockyou} (\texttt{ma*}) com regras \emph{best64}/\emph{leet}/anos, 3 \emph{encodings} & $\sim$15 funções de \emph{hash} & falhou\\
\emph{Passphrases}/frases temáticas (Enigma/cripto; PL/DE/EN/PT; 1--3 palavras) & milhões de combinações & falhou\\
Nomes do grupo e dos criptoanalistas do Biuro Szyfrów & milhares de combinações & falhou\\
Construções alternativas (\emph{double-hash}; entrada hex/base64/revertida; UTF-16/32) & dezenas $\times$ 7--12 \emph{hashes} & falhou\\
``Textbook RSA'' na chave da mensagem (expoente pequeno; contornar a senha) & $e\in\{3,\dots,65537\}$, big/little-endian & falhou\\
Fatoração elementar de $N$ (Fermat, Pollard $p\!-\!1$, Brent-$\rho$; Python) & válida só p/ primos fracos & falhou\\
\midrule
\textbf{Combinador de 2 palavras} (palavra-\texttt{ma} $\times$ dicionário amplo; C+OpenMP) & SHA-256 & \textbf{sucesso}\\
\bottomrule
\end{tabular}
\caption{Resumo das abordagens testadas na quebra da senha e na fatoração de $N$.}
\label{tab:resumo}
\end{table}

\section{Reprodutibilidade --- passo a passo}
\begin{enumerate}[leftmargin=*,nosep]
  \item \textbf{Senha:} compilar e rodar o combinador\\
  \texttt{gcc -O3 -fopenmp -o brute\_combine brute\_combine.c}\\
  \texttt{./brute\_combine mawords.txt allwords.txt 10} $\Rightarrow$ \texttt{marvinminsky}.
  \item \textbf{Chave pública:} \texttt{openssl enc -d -aes-256-ctr -nosalt -md sha256 -pass pass:marvinminsky -in key\_public.en}.
  \item \textbf{Extrair $N,e$:} \texttt{openssl rsa -pubin -in pub.pem -text -noout}.
  \item \textbf{Fatorar $N$:} \texttt{./cado-nfs.py <N> -t 4} $\Rightarrow p,q$.
  \item \textbf{Mensagem:} \texttt{python3 finish.py <p> <q>} (reconstrói a privada, decifra a chave
  simétrica e a mensagem).
\end{enumerate}

\section{Código-fonte}
\subsection{\texttt{brute\_combine.c} --- combinador de 2 palavras (SHA-256, C+OpenMP)}
\lstinputlisting[language=C]{../src/brute_combine.c}

\subsection{\texttt{decrypt\_pubkey.py} --- decifração da chave pública}
\lstinputlisting[language=Python]{../src/decrypt_pubkey.py}

\subsection{\texttt{factor\_rsa.py} --- fatoração elementar (pré-GNFS)}
\lstinputlisting[language=Python]{../src/factor_rsa.py}

\subsection{\texttt{finish.py} --- reconstrução da chave privada e decifração final}
\lstinputlisting[language=Python]{../src/finish.py}

\section*{Conclusão}
\addcontentsline{toc}{section}{Conclusão}
A cadeia completa foi quebrada: \texttt{senha=marvinminsky} $\rightarrow$ chave pública RSA-382
$\rightarrow$ fatoração por GNFS $\rightarrow$ \texttt{chave=shafigoldwas} $\rightarrow$ mensagem em
claro endereçada ao Grupo~2. O exercício ilustra na prática por que \textbf{RSA de baixa entropia}
(módulos pequenos e fatoráveis) e \textbf{senhas de espaço de busca limitado} (nomes conhecidos) são
inseguros, mesmo combinados com AES-256.

\end{document}
